\documentclass[aspectratio=169,11pt]{beamer}
\usetheme{Madrid}
\usecolortheme{default}
\setbeamertemplate{navigation symbols}{}
\setbeamertemplate{footline}[frame number]
\usepackage[utf8]{inputenc}
\usepackage[T1]{fontenc}
\usepackage{lmodern}
\usepackage{amsmath}
\usepackage{booktabs}
\usepackage{array}
\usepackage{ragged2e}
\usepackage{xcolor}
\usepackage{pgfpages}

\definecolor{UNBCGreen}{RGB}{0,102,51}
\definecolor{SoftGold}{RGB}{189,155,96}
\setbeamercolor{structure}{fg=UNBCGreen}
\setbeamercolor{title}{fg=UNBCGreen}
\setbeamercolor{frametitle}{fg=UNBCGreen,bg=UNBCGreen!8}
\setbeamercolor{block title}{fg=white,bg=UNBCGreen}
\setbeamercolor{block body}{bg=UNBCGreen!5}

\title{Economics Is Everywhere}
\subtitle{A 20-minute introduction for students and parents}
\author{Prepared from the talk draft}
\date{}
\setbeameroption{show notes on second screen=right}

\begin{document}

\begin{frame}
  \titlepage
  \note{Welcome the audience and briefly introduce yourself. Say this talk is designed for students and parents who may be new to economics. Emphasize that the goal is to show that economics is practical, interesting, and relevant to real life. Mention that the presentation will be light, accessible, and connected to everyday examples.}
\end{frame}

\begin{frame}{What will we cover?}
\begin{itemize}
    \item What economics really is
    \item Why it matters in everyday life
    \item Big issues economists study
    \item Fascinating fields like behavioural economics and neuroeconomics
    \item Why students should consider studying economics
\end{itemize}
\vspace{0.4cm}
\textit{Short version: economics is not just about money - it is about choices, incentives, and real life.}
\note{Walk the audience through the roadmap. Tell them you will start with a simple definition of economics, then connect it to daily life, then move to some exciting fields within economics, and finally explain why economics is a useful field of study. This is also a good place to lower anxiety by saying there will be no complicated math in this talk.}
\end{frame}

\begin{frame}{What is economics?}
\begin{block}{Simple idea}
Economics is the study of how people, businesses, and governments make choices when resources are limited.
\end{block}
\vspace{0.2cm}
\begin{itemize}
    \item We all face scarcity: limited time, money, energy, and opportunities
    \item Because resources are limited, every choice involves a \textbf{trade-off}
    \item Economics helps us think clearly about those trade-offs
\end{itemize}
\vspace{0.3cm}
\textbf{Example:} A student wants sleep, good grades, a part-time job, friends, and time to relax.\newline
There are only 24 hours in a day. That is economics.
\note{Define scarcity in very simple language. Explain that economics begins with the fact that we cannot have everything at the same time, so choices must be made. Use the student example and invite the audience to smile at how familiar it sounds. You can add that parents face the same issue with time, budgets, and priorities. Stress that economics is really about decision-making under limits.}
\end{frame}

\begin{frame}{Economics is everywhere}
\begin{columns}[T,totalwidth=\textwidth]
\column{0.5\textwidth}
\textbf{At home}
\begin{itemize}
    \item Grocery budgets
    \item Saving vs spending
    \item Rent, mortgage, and bills
\end{itemize}
\vspace{0.2cm}
\textbf{At school}
\begin{itemize}
    \item Choosing courses
    \item Managing time
    \item Deciding whether to work while studying
\end{itemize}

\column{0.5\textwidth}
\textbf{In society}
\begin{itemize}
    \item Inflation and wages
    \item Housing affordability
    \item Taxes and public services
    \item Climate policy
\end{itemize}
\end{columns}
\vspace{0.4cm}
\textit{Economics is what you are already doing - even if you never call it economics.}
\note{Use this slide to connect quickly with both students and parents. Mention that whenever a family makes a budget, a student chooses between work and study time, or society debates housing and taxes, economics is present. Emphasize that economics is not remote or abstract. It is already part of daily life.}
\end{frame}

\begin{frame}{A few real-life examples}
\begin{itemize}
    \item \textbf{Coffee shop line:} high demand + limited supply = waiting time
    \item \textbf{Concert tickets:} scarce seats + huge demand = higher resale prices
    \item \textbf{Grocery shopping:} families substitute products when prices rise
    \item \textbf{Choosing a degree:} every decision has an \textbf{Opportunity cost}
\end{itemize}
\vspace{0.4cm}
\begin{block}{Opportunity cost}
The real cost of a choice is the value of the next best alternative you give up.
\end{block}
\vspace{0.2cm}
\textit{If you spend 3 hours scrolling on your phone instead of studying, economics would politely ask: ``What did that really cost you?''}
\note{Slow down here and explain opportunity cost carefully because it is one of the most important ideas in economics. Give one student example and one family example. You can joke that opportunity cost is why late-night scrolling often becomes expensive in the morning. This slide works well when spoken conversationally.}
\end{frame}

\begin{frame}{The big questions economics helps explain}
\begin{itemize}
    \item Why do prices rise? \hfill \textbf{Inflation}
    \item Why can jobs be hard to find? \hfill \textbf{Unemployment}
    \item Why is housing expensive? \hfill \textbf{Supply, demand, policy}
    \item Why do some people or places do better than others? \hfill \textbf{Growth and inequality}
    \item How should we respond to pollution and climate change? \hfill \textbf{Environmental economics}
\end{itemize}
\vspace{0.5cm}
\textbf{Key lesson:} economics teaches us to look beyond the first effect and ask,\newline
\textit{``What happens next?''}
\note{Use this slide to show that economics helps explain the big issues people hear about in the news. Keep the explanations intuitive. For inflation, mention groceries and rent. For unemployment, mention job opportunities for young people. For housing, mention affordability in many Canadian communities. Emphasize that economics is useful because it encourages us to think past the immediate effect and consider consequences over time.}
\end{frame}

\begin{frame}{Economics is about people, not just numbers}
\begin{itemize}
    \item Families deciding how to spend their income
    \item Students choosing their future
    \item Workers looking for secure jobs
    \item Governments deciding how to use limited budgets
    \item Communities trying to improve quality of life
\end{itemize}
\vspace{0.5cm}
\begin{block}{Important point}
Good economics is deeply human. It studies well-being, fairness, incentives, and the consequences of choices.
\end{block}
\note{This slide is important for correcting the stereotype that economics is cold or mechanical. Tell the audience that behind every statistic are real people and real lives. Mention that economics studies jobs, family budgets, access to housing, and opportunities for communities. This helps humanize the discipline.}
\end{frame}

\begin{frame}{Interesting field 1: Behavioural economics}
\begin{itemize}
    \item Traditional economics often assumes people are rational
    \item Behavioural economics studies how people \textbf{actually} behave
    \item It examines procrastination, habits, overconfidence, peer effects, and self-control problems
\end{itemize}
\vspace{0.3cm}
\textbf{Examples}
\begin{itemize}
    \item Why do we buy things we do not need?
    \item Why do people delay saving for retirement?
    \item Why is ``starting tomorrow'' such a popular life strategy?
\end{itemize}
\vspace{0.3cm}
\textit{Behavioural economics reminds us that humans are not robots. Sometimes we are not even very good spreadsheets.}
\note{This is a fun slide. Use relatable examples like procrastinating on homework, buying things during online sales, or saying 'I will start tomorrow.' Explain that behavioural economics studies the gap between what people plan to do and what they actually do. This is often one of the most memorable parts of the talk for students.}
\end{frame}

\begin{frame}{Interesting field 2: Neuroeconomics}
\begin{itemize}
    \item Neuroeconomics connects economics, psychology, and neuroscience
    \item It studies what happens in the brain when people make decisions
    \item It helps explain risk-taking, rewards, trust, addiction, and impulsive choices
\end{itemize}
\vspace{0.4cm}
\textbf{Questions it asks}
\begin{itemize}
    \item Why do people take risks?
    \item How do emotions affect financial decisions?
    \item Why does immediate reward often beat long-term benefit?
\end{itemize}
\vspace{0.3cm}
\textit{No, it does not read minds. Economics still struggles to predict what a family wants for dinner when everyone says, ``Anything is fine.''}
\note{Introduce neuroeconomics as an example of how modern economics connects with other sciences. Keep it simple: this field looks at how the brain is involved in decision-making. Mention that it helps explain risk, reward, trust, and impulsive choices. Use the dinner joke to keep the tone light.}
\end{frame}

\begin{frame}{Other fascinating fields in economics}
\begin{columns}[T,totalwidth=\textwidth]
\column{0.5\textwidth}
\begin{itemize}
    \item \textbf{Labour economics}\newline jobs, wages, skills, unemployment
    \item \textbf{Health economics}\newline healthcare costs and policy
    \item \textbf{Environmental economics}\newline climate, pollution, sustainability
\end{itemize}
\column{0.5\textwidth}
\begin{itemize}
    \item \textbf{Development economics}\newline poverty, education, growth
    \item \textbf{Financial economics}\newline banking, investing, risk
    \item \textbf{Public economics}\newline taxes, spending, government policy
\end{itemize}
\end{columns}
\vspace{0.4cm}
\textbf{Bottom line:} economics is broad, modern, and connected to major issues in the world.
\note{Move briskly through this slide and show the range of the discipline. You do not need to define every field in detail. The goal is to let the audience see that economics is much broader than they may have thought. Mention that students can find areas that connect to jobs, health, sustainability, public policy, finance, and global development.}
\end{frame}

\begin{frame}{Why economics matters now more than ever}
\begin{itemize}
    \item Rising living costs and inflation
    \item Housing and affordability pressures
    \item Artificial intelligence and changing labour markets
    \item Global trade, supply chains, and uncertainty
    \item Climate and sustainability challenges
\end{itemize}
\vspace{0.4cm}
\begin{block}{Why study it?}
Economics gives students a framework for understanding today's world instead of just reacting to headlines.
\end{block}
\note{This slide connects economics to the world students and parents are living in right now. Use current examples like rising prices, housing pressure, and technology changing work. Emphasize that economics helps people interpret the world more thoughtfully instead of simply reacting to news stories.}
\end{frame}

\begin{frame}{What skills does economics build?}
\begin{itemize}
    \item Critical thinking
    \item Problem solving
    \item Data interpretation
    \item Decision-making under uncertainty
    \item Clear writing and communication
    \item Understanding incentives and unintended consequences
\end{itemize}
\vspace{0.4cm}
\textit{These are not only economist skills - they are career skills and life skills.}
\note{Explain that economics is valuable even for students who do not plan to become economists. The field develops analytical and transferable skills that employers value across many sectors. You can mention that economics trains people to organize evidence, think clearly, and communicate arguments well.}
\end{frame}

\begin{frame}{Why students should consider economics}
\begin{itemize}
    \item It is practical and relevant to everyday life
    \item It combines numbers, ideas, and real-world problems
    \item It works well on its own or with business, political science, data science, environmental studies, and more
    \item It opens doors to many careers
\end{itemize}
\vspace{0.4cm}
\textbf{Common career paths}
\begin{itemize}
    \item Government and public policy
    \item Banking and finance
    \item Business and consulting
    \item Research and analytics
    \item Law, education, and graduate study
\end{itemize}
\note{This is where you become more persuasive. Stress that economics is flexible and opens many pathways. Speak clearly about jobs and employability because that matters to both students and parents. You can note that economics works very well as a major, a minor, or as part of a combined skill set with another discipline.}
\end{frame}

\begin{frame}{A message for parents and students}
\begin{block}{For students}
If you are curious about how the world works and enjoy asking ``why,'' economics is worth exploring.
\end{block}
\vspace{0.3cm}
\begin{block}{For parents}
Economics develops adaptable, analytical, employable graduates with skills valued across many sectors.
\end{block}
\vspace{0.4cm}
\textit{Also, after one economics course, students may finally understand why ``we cannot buy everything'' is not just parenting - it is scarcity.}
\note{Address both groups directly. Tell students that curiosity matters more than already knowing economics. Tell parents that economics develops practical and adaptable graduates. End the slide with the scarcity joke to keep the tone friendly and memorable.}
\end{frame}

\begin{frame}{Final takeaway}
\begin{block}{Economics is everywhere}
It helps us understand choices, incentives, markets, policy, and human behaviour.
\end{block}
\vspace{0.4cm}
\begin{itemize}
    \item It is relevant to everyday life
    \item It helps explain major challenges in society
    \item It includes exciting fields like behavioural economics and neuroeconomics
    \item It builds valuable skills and strong career opportunities
\end{itemize}
\vspace{0.5cm}
\centering
\Large \textbf{Economics is not just a subject. It is a way of understanding the world.}
\note{Use this slide to summarize confidently. Reiterate that economics is practical, broad, and useful. Return to the central message that economics helps people understand the world around them. Keep your voice slightly slower here so the talk feels like it is concluding clearly.}
\end{frame}

\begin{frame}[plain]
\centering
\vfill
{\Huge Thank you!}\par
\vspace{0.3cm}
Questions?
\vfill
\note{Thank the audience warmly and invite questions. If there is time, encourage students and parents to ask about programs, careers, or specific areas of economics that interested them during the talk.}
\end{frame}

\end{document}
